\heading{经济学原理（II）-模拟题3-期中}
\noteinfo{这套题比前两套更难，重点加强“口径判断 + 多步计算 + 模型比较”。不少题目并非难在公式，而是难在先判断该用哪个口径。}

\instructions{建议答题时间 120 分钟。请特别注意：本卷多处要求你解释“为什么”，只写结论通常不给满分。}

\textbf{一、单项选择题}

\begin{question}
某国生产可能性边界外移，但其生产点仍留在边界内侧。这最恰当的解释是
\begin{choiceoptions}[2]
  \optionitem{该国一定实现了帕累托最优}
  \optionitem{资源配置能力可能改善了，但当前仍存在闲置资源}
  \optionitem{机会成本一定下降到零}
  \optionitem{该国已不再面临稀缺性}
\end{choiceoptions}
\begin{solution}
边界外移表示生产能力提升，但点在边界内侧说明现实生产没有用满资源，仍存在无效率。\anstext{B}
\end{solution}
\end{question}

\begin{question}
下列哪一项会计入 2026 年中国 GDP，\textbf{但不计入} 2026 年中国 GNP？
\begin{choiceoptions}[1]
  \optionitem{中国居民持有美国股票获得的股息}
  \optionitem{在中国境内生产、由外资企业实现并最终汇回国外母公司的新增产出}
  \optionitem{中国政府向大学生发放生活补助}
  \optionitem{中国居民购买二手房支付给卖方的房款}
\end{choiceoptions}
\begin{solution}
GDP 是地域口径，GNP 是居民口径。境内外资企业在中国生产的新增产出计入 GDP，但收益归外国居民，不增加中国 GNP。\anstext{B}
\end{solution}
\end{question}

\begin{question}
若某商品是进口消费品且价格大幅上涨，在其他条件不变时，这一冲击最直接导致
\begin{choiceoptions}[2]
  \optionitem{CPI 上升，而 GDP 平减指数不一定同方向变化}
  \optionitem{GDP 平减指数上升，而 CPI 必定不变}
  \optionitem{CPI 与 GDP 平减指数必定同比例上升}
  \optionitem{真实 GDP 必定下降}
\end{choiceoptions}
\begin{solution}
CPI 包含居民消费的进口品，而 GDP 平减指数只覆盖国内生产。进口品涨价会直接推高 CPI，但不一定直接推高 GDP 平减指数。\anstext{A}
\end{solution}
\end{question}

\begin{question}
在买方垄断劳动市场中，若政府设置一个适度的最低工资，可能提高就业，最关键的机制是
\begin{choiceoptions}[2]
  \optionitem{最低工资让劳动供给曲线完全水平}
  \optionitem{最低工资降低了工人的边际产品}
  \optionitem{最低工资可能使边际雇佣成本下降到更接近工资本身}
  \optionitem{最低工资必然使企业利润上升}
\end{choiceoptions}
\begin{solution}
买方垄断下边际雇佣成本高于工资。适度最低工资会让企业在一段区间内以相同工资雇更多工人，从而降低边际雇佣成本。\anstext{C}
\end{solution}
\end{question}

\begin{question}
若某变量每 14 年翻一倍，根据 70 规则，其年均增长率大约是
\begin{choiceoptions}[2]
  \optionitem{3\%}
  \optionitem{5\%}
  \optionitem{7\%}
  \optionitem{14\%}
\end{choiceoptions}
\begin{solution}
按 70 规则，增长率约为 $70/14=5\%$。\anstext{B}
\end{solution}
\end{question}

\begin{question}
在索罗模型中，若实际人均投资持续小于 $(n+\delta)k$，则长期上意味着
\begin{choiceoptions}[2]
  \optionitem{资本密度上升并趋向更高稳态}
  \optionitem{资本密度下降，经济位于稳态右侧}
  \optionitem{资本密度不变，已在稳态}
  \optionitem{人口增长率为零}
\end{choiceoptions}
\begin{solution}
若 $sf(k)<(n+\delta)k$，人均资本被折旧与人口增长“摊薄”得更多，故 $k$ 会下降，说明经济位于稳态资本密度的右侧。\anstext{B}
\end{solution}
\end{question}

\begin{question}
关于预算赤字与可贷资金市场，下列说法最准确的是
\begin{choiceoptions}[1]
  \optionitem{预算赤字只影响公共储蓄，不影响国民储蓄}
  \optionitem{预算赤字会使可贷资金供给右移，真实利率下降}
  \optionitem{预算赤字会降低公共储蓄，从而减少国民储蓄并挤出投资}
  \optionitem{预算赤字只影响名义利率，不影响真实利率}
\end{choiceoptions}
\begin{solution}
赤字意味着 $T-G$ 下降，公共储蓄减少，国民储蓄减少，可贷资金供给左移，真实利率上升，投资被挤出。\anstext{C}
\end{solution}
\end{question}

\begin{question}
若两种资产的个别风险都很大，但二者收益率相关系数为 $-1$，则在适当权重下
\begin{choiceoptions}[2]
  \optionitem{组合风险一定大于个别资产风险}
  \optionitem{组合风险可能被完全消除}
  \optionitem{系统性风险一定上升}
  \optionitem{期望收益一定为零}
\end{choiceoptions}
\begin{solution}
完全负相关允许在合适权重下构造无风险组合。\anstext{B}
\end{solution}
\end{question}

\begin{question}
关于有效市场假说，下列说法最恰当的是
\begin{choiceoptions}[1]
  \optionitem{若市场有效，则资产价格永远不会偏离基本面}
  \optionitem{半强式有效意味着公开信息已迅速反映在价格中}
  \optionitem{有效市场假说说明内幕信息也不可能带来超额收益}
  \optionitem{只要价格波动很大，就一定违背有效市场假说}
\end{choiceoptions}
\begin{solution}
半强式有效指公开信息迅速进入价格；它不等于价格绝不会短时偏离基本面。\anstext{B}
\end{solution}
\end{question}

\begin{question}
在两期消费模型中，若真实利率 $r$ 上升而偏好不变，则通常直接改变的是
\begin{choiceoptions}[2]
  \optionitem{预算线斜率}
  \optionitem{总禀赋现值}
  \optionitem{基尼系数}
  \optionitem{资本份额 $\alpha$
}
\end{choiceoptions}
\begin{solution}
真实利率改变现期与未来消费的相对价格，因此改变跨期预算线斜率。\anstext{A}
\end{solution}
\end{question}

\textbf{二、简答题}

\begin{question}
比较“支付能力原则”和“受益原则”。为什么税制讨论中，纵向公平和横向公平更常与支付能力原则联系在一起？
\begin{solution}
受益原则主张谁从公共服务中受益更多，谁承担更多税负；支付能力原则则主张按纳税人的承受能力征税。纵向公平与横向公平都在讨论“具有不同或相近支付能力的人应承担怎样的税负”，因此更自然地从支付能力原则引申出来。受益原则更适合解释某些特定收费或准价格机制，但在普遍税制里，公共品受益难以精确度量，所以实际讨论更常围绕支付能力展开。\end{solution}
\end{question}

\begin{question}
请比较 CPI 与 GDP 平减指数：二者为什么会在“进口能源价格大涨、但国内政府购买品价格稳定”的情形下给出不同通胀信号？
\begin{solution}
CPI 衡量代表性消费者消费篮子的成本，进口能源若进入居民消费，会直接推高 CPI。GDP 平减指数只覆盖国内生产的最终产品与服务，不直接把进口品涨价计入；相反，政府购买品若价格稳定，会使平减指数受到的上推压力较小。因此二者的差异来自\textbf{覆盖范围不同}：一个是消费口径，一个是国内生产口径。\end{solution}
\end{question}

\begin{question}
“贫穷国家未必一定追赶上富裕国家。”请用索罗模型中的“条件收敛”思想解释这句话。
\begin{solution}
索罗模型并不是说所有贫穷国家都会自动追上富裕国家，而是说在\textbf{储蓄率、人口增长、折旧、技术水平和制度环境相同}的条件下，资本较少的经济体因为资本边际产出更高，会更快增长并向同一稳态收敛。若各国参数不同，例如储蓄率低、人口增长高、制度差、技术落后，则它们会收敛到不同稳态，贫穷国家就可能长期仍然贫穷。因此收敛是有条件的，而不是无条件的。\end{solution}
\end{question}

\begin{question}
请解释系统性风险与非系统性风险的区别，并说明为什么金融危机时期“分散化失灵”的说法并不意味着分散化原理本身是错的。
\begin{solution}
非系统性风险是某家公司、某个行业特有的风险，例如管理失误、产品召回；系统性风险是影响整个市场或宏观经济的风险，例如金融危机、普遍衰退、通胀冲击。分散化能显著降低前者，却无法消除后者。危机时期大量资产一起下跌，是因为共同暴露在系统性冲击下，相关性上升，所以分散化效果减弱。但这并不推翻分散化原理；它只是说明分散化主要对付的是特有风险，而不是所有风险。\end{solution}
\end{question}

\textbf{三、计算和分析题}

\begin{question}
某经济体有三家企业，企业 3 的全部产品卖给家庭。请补全下表，并用至少两种方法验证 GDP。
\begin{center}
\begin{tabular}{c|cccc}
 & 企业1 & 企业2 & 企业3 & 经济整体 \\\hline
销售收入 & 800 & 1500 & \blank[1.2cm] & -- \\
原料投入 & 0 & \blank[1.2cm] & 1500 & -- \\
劳动报酬 & 500 & 250 & 700 & \blank[1.2cm] \\
地租 & 100 & 100 & 200 & \blank[1.2cm] \\
利息 & 50 & 50 & 100 & \blank[1.2cm] \\
利润 & 150 & 300 & 400 & \blank[1.2cm] \\
增加值 & \blank[1.2cm] & \blank[1.2cm] & 1400 & \blank[1.2cm]
\end{tabular}
\end{center}
\begin{solution}
企业 1 的增加值是 $800-0=800$。企业 2 的原料投入来自企业 1，因此是 800，增加值为
\[
1500-800=700.
\]
企业 3 销售收入为
\[
1500+1400=2900.
\]
劳动报酬总额为 $500+250+700=1450$，地租总额为 $400$，利息总额为 $200$，利润总额为 $850$。GDP 等于
\[
800+700+1400=2900.
\]
收入法验证：
\[
1450+400+200+850=2900.
\]
支出法验证：企业 3 的最终产品全部卖给家庭，因此最终支出也为 2900。\ansmath{GDP=2900}
\end{solution}
\end{question}

\begin{question}
某国五等份收入份额为：4\%、9\%、15\%、24\%、48\%。
\begin{enumerate}[label=(\alph*),leftmargin=2.5em]
  \item 写出洛伦兹曲线的关键点；
  \item 用折线近似计算基尼系数；
  \item 说明只用五等份数据估计的基尼系数，通常为何会低估真实不平等程度。
\end{enumerate}
\begin{solution}
累积收入份额依次为
\[
0,\ 0.04,\ 0.13,\ 0.28,\ 0.52,\ 1.
\]
因此关键点为
\[
(0,0),(0.2,0.04),(0.4,0.13),(0.6,0.28),(0.8,0.52),(1,1).
\]
洛伦兹曲线下方面积为各梯形面积之和：
\[
A_L=0.2\times\frac{0+0.04}{2}+0.2\times\frac{0.04+0.13}{2}+0.2\times\frac{0.13+0.28}{2}+0.2\times\frac{0.28+0.52}{2}+0.2\times\frac{0.52+1}{2}.
\]
计算得
\[
A_L=0.294.
\]
故基尼系数
\[
G=1-2A_L=1-0.588=0.412.
\]
\ansmath{G\approx 0.412}

只用五等份时，我们把每一组内部收入看成“均匀分布”，会抹平组内差异，尤其会低估最高收入组内部的厚尾集中，因此通常会低估真实不平等。\end{solution}
\end{question}

\begin{question}
一个经济体生产甲、乙、丙三种产品，数据如下：
\begin{center}
\begin{tabular}{c|cc|cc|cc}
 & \multicolumn{2}{c|}{2024年} & \multicolumn{2}{c|}{2025年} & \multicolumn{2}{c}{2026年} \\
商品 & 产量 & 价格 & 产量 & 价格 & 产量 & 价格 \\\hline
甲 & 20 & 2 & 25 & 3 & 24 & 4 \\
乙 & 30 & 1 & 35 & 1 & 40 & 2 \\
丙 & 10 & 5 & 12 & 6 & 15 & 6
\end{tabular}
\end{center}
\begin{enumerate}[label=(\alph*),leftmargin=2.5em]
  \item 计算 2025 至 2026 年名义 GDP 增长率；
  \item 以 2025 年为基期，求 2024 年真实 GDP 与 2026 年真实 GDP；
  \item 计算 2026 年 GDP 平减指数（以 2025 年为基期）；
  \item 若代表性消费者的固定篮子为甲 2、乙 3、丙 1，以 2025 年为基期，求 2026 年 CPI，并解释其与 GDP 平减指数为何不同。
\end{enumerate}
\begin{solution}
2025 年名义 GDP：
\[
25\times3+35\times1+12\times6=182.
\]
2026 年名义 GDP：
\[
24\times4+40\times2+15\times6=266.
\]
名义 GDP 增长率：
\[
\frac{266-182}{182}\times100\%=46.15\%.
\]

以 2025 年价格计算：
\[
RGDP_{2024|2025}=20\times3+30\times1+10\times6=150,
\]
\[
RGDP_{2026|2025}=24\times3+40\times1+15\times6=202.
\]

GDP 平减指数：
\[
P^{GDP}_{2026|2025}=\frac{266}{202}\times100\approx 131.68.
\]

固定篮子在 2025 年的成本：
\[
2\times3+3\times1+1\times6=15.
\]
2026 年成本：
\[
2\times4+3\times2+1\times6=20.
\]
因此
\[
CPI_{2026|2025}=\frac{20}{15}\times100=133.33.
\]
CPI 与平减指数不同，是因为 CPI 使用固定消费篮子，而 GDP 平减指数使用当期国内生产构成，权重不同。\ansmath{\text{名义 GDP 增长率}\approx 46.15\%,\ RGDP_{2024|2025}=150,\ RGDP_{2026|2025}=202}
\ansmath{P^{GDP}_{2026|2025}\approx 131.68,\ CPI_{2026|2025}=133.33}
\end{solution}
\end{question}

\begin{question}
某经济体的人均生产函数为
\[
y=k^{1/3},
\]
储蓄率 $s=0.25$，人口增长率 $n=0.02$，折旧率 $\delta=0.04$。回答：
\begin{enumerate}[label=(\alph*),leftmargin=2.5em]
  \item 求稳态资本密度 $k^*$ 与稳态人均产出 $y^*$；
  \item 若当前 $k=100$，判断该经济体将向上还是向下调整，并说明原因；
  \item 求黄金规则储蓄率，并判断现行储蓄率高于还是低于黄金规则。
\end{enumerate}
\begin{solution}
稳态满足
\[
0.25k^{1/3}=0.06k.
\]
移项得
\[
k^{2/3}=\frac{0.25}{0.06}=\frac{25}{6}.
\]
因此
\[
k^*=\left(\frac{25}{6}\right)^{3/2}\approx 8.505.
\]
稳态人均产出为
\[
y^*=(k^*)^{1/3}=\left(\frac{25}{6}\right)^{1/2}\approx 2.041.
\]

若当前 $k=100$，则
\[
sf(k)=0.25\times 100^{1/3}\approx 1.160,
\]
\[
(n+\delta)k=0.06\times 100=6.
\]
因前者小于后者，资本密度会下降，所以经济将向下调整，说明当前位于稳态右侧。

Cobb--Douglas 情形下，黄金规则储蓄率等于资本份额 $\alpha=1/3$。由于现行储蓄率 $0.25<1/3$，低于黄金规则。\ansmath{k^*\approx 8.505,\quad y^*\approx 2.041,\quad s_{gold}=\frac13}
\end{solution}
\end{question}

\begin{question}
考虑可贷资金市场。起初政府预算平衡，随后发生两件事：
\begin{enumerate}[label=(\arabic*),leftmargin=2.5em]
  \item 政府出现预算赤字；
  \item 同时企业获得投资税收抵免，投资需求上升。
\end{enumerate}
请回答：
\begin{enumerate}[label=(\alph*),leftmargin=2.5em]
  \item 两件事分别使哪条曲线移动、向哪个方向移动？
  \item 对均衡真实利率的影响是否确定？对均衡投资量是否确定？
  \item 若预期通胀率同时上升 3 个百分点，而真实利率长期不变，名义利率如何变化？
\end{enumerate}
\begin{solution}
预算赤字降低公共储蓄，使可贷资金\textbf{供给曲线左移}；投资税收抵免提高投资意愿，使\textbf{投资需求曲线右移}。

两者都会推高均衡真实利率，因此真实利率\textbf{确定上升}。但均衡投资量的变化不确定：需求右移会提高投资，供给左移会压低投资，净效应取决于哪一个力量更强。

若长期真实利率不变，预期通胀率上升 3 个百分点，则按费雪方程
\[
i\approx r+\pi^e,
\]
名义利率也上升约 3 个百分点。\end{solution}
\end{question}

\begin{question}
两种资产 $A$、$B$ 在两种状态下的收益率如下，概率各为 0.5：
\begin{center}
\begin{tabular}{c|cc}
状态 & $A$ & $B$\\\hline
1 & 30\% & 0\%\\
2 & -10\% & 20\%
\end{tabular}
\end{center}
\begin{enumerate}[label=(\alph*),leftmargin=2.5em]
  \item 分别求 $A$ 与 $B$ 的期望收益率和标准差；
  \item 若投资者用权重 $w$ 配置给 $A$，其余配置给 $B$，写出组合在两种状态下的收益率；
  \item 求使两种状态收益率相等的权重 $w$，并给出此时无风险收益率；
  \item 解释为什么这个例子说明“相关性而非单个方差”决定分散化效果。
\end{enumerate}
\begin{solution}
\[
E(r_A)=0.5\times30\%+0.5\times(-10\%)=10\%,
\]
\[
E(r_B)=0.5\times0\%+0.5\times20\%=10\%.
\]
标准差：
\[
\sigma_A=\sqrt{0.5(20\%)^2+0.5(-20\%)^2}=20\%,
\]
\[
\sigma_B=\sqrt{0.5(-10\%)^2+0.5(10\%)^2}=10\%.
\]

组合在状态 1 的收益率：
\[
r_1=0.3w.
\]
状态 2：
\[
r_2=-0.1w+0.2(1-w)=0.2-0.3w.
\]
令二者相等：
\[
0.3w=0.2-0.3w\Rightarrow w=\frac13.
\]
此时无风险收益率为
\[
r_f=0.3\times\frac13=10\%.
\]

该例中，单个资产方差并不小，但由于两种状态下的波动方向相反，恰当配置后可以抵消；因此分散化效果取决于资产之间如何共同波动，而不只取决于每个资产“自己有多波动”。\ansmath{E(r_A)=E(r_B)=10\%,\ \sigma_A=20\%,\ \sigma_B=10\%,\ w=\frac13,\ r_f=10\%}
\end{solution}
\end{question}

\begin{question}
两期消费模型中，消费者效用为
\[
u(c_1,c_2)=\ln c_1+0.9\ln c_2,
\]
初始资产为 0，真实利率为 $r$，劳动收入分别为 $y_1,y_2$。
\begin{enumerate}[label=(\alph*),leftmargin=2.5em]
  \item 写出欧拉方程；
  \item 写出终身预算约束；
  \item 解出最优 $c_1$，并求其对 $y_1$ 的边际消费倾向 MPC。
\end{enumerate}
\begin{solution}
欧拉方程：
\[
\frac{1}{c_1}=0.9(1+r)\frac{1}{c_2}
\Rightarrow c_2=0.9(1+r)c_1.
\]
终身预算约束：
\[
c_1+\frac{c_2}{1+r}=y_1+\frac{y_2}{1+r}.
\]
代入得
\[
c_1+\frac{0.9(1+r)c_1}{1+r}=1.9c_1=y_1+\frac{y_2}{1+r}.
\]
故
\[
c_1=\frac{1}{1.9}\left(y_1+\frac{y_2}{1+r}\right).
\]
对 $y_1$ 求导得
\[
MPC=\frac{\partial c_1}{\partial y_1}=\frac{1}{1.9}=\frac{10}{19}.
\]
\ansmath{c_1=\frac{1}{1.9}\left(y_1+\frac{y_2}{1+r}\right),\quad MPC=\frac{10}{19}}
\end{solution}
\end{question}

